
\begin{frame}{Definition 1: Demographic Parity (Independence)}

    \begin{block}{Demographic parity / statistical parity}
        \[
            \Pr(\Yhat = 1 \mid A = \GrpA) \;=\; \Pr(\Yhat = 1 \mid A = \GrpB)
        \]
        The model selects the same \emph{proportion} of each group.
    \end{block}

    \vspace{0.6em}

    \begin{columns}[T]
        \begin{column}{0.48\textwidth}
            \textbf{When it is the right question}
            \begin{itemize}\small
                \setlength\itemsep{0.3em}
                \item Allocating an opportunity where you do not trust the label
                      (advertising, outreach, shortlisting).
                \item You believe the historical outcome $Y$ is itself contaminated.
                \item A widely used convention: a selection-rate ratio below $0.8$ between groups
                      is treated as a gap worth investigating. Most audit tooling flags it by
                      default.
            \end{itemize}
        \end{column}
        \begin{column}{0.48\textwidth}
            \textbf{Where it goes wrong}
            \begin{itemize}\small
                \setlength\itemsep{0.3em}
                \item It ignores $Y$ entirely. A model can satisfy it by selecting qualified people
                      from one group and arbitrary people from the other.
                \item If the true base rates genuinely differ, enforcing it forces errors.
                \item Measured as a \textbf{ratio} (disparate impact) or a \textbf{difference}
                      --- report which.
            \end{itemize}
        \end{column}
    \end{columns}
\end{frame}
